\section{Data}
In this study, we use data products generated by the weekly intermittent Data Reduction Pipelines (DRPs). 
These DRPs aim to test current data quality and produce LSST-like downstream data from the single-visit exposures. 
We use two main data products: the \texttt{object} and the \texttt{refit\_psf\_star} tables.

The \texttt{object} table contains data of stars and galaxies derived from the per-band coadded images. 
The \texttt{refit\_psf\_star} table contains data of stars and galaxies derived from the single-visit images and PSF moments modelled with the \texttt{PiFF} package \citep{jarvis2021}.

For measured second-order moments, we use \texttt{slot\_Shape\_xx}, \texttt{slot\_Shape\_yy} and \texttt{slot\_Shape\_xy} that have been measured using the HSM method \citep{2003MNRAS.343..459H,2005MNRAS.361.1287M, 2018MNRAS.481.3170M} in which the moments are determined using an iterative elliptical Gaussian fit.
For the second-order moments of the \texttt{PiFF} PSF model, we use \texttt{slot\_PsfShape\_xx}, \texttt{slot\_PsfShape\_yy} and \texttt{slot\_PsfShape\_xy}, also derived using the HSM shape algorithm.

For star magnitudes, we use calibrated magnitudes (\texttt{ap12Flux\_mag}) from the \texttt{refit\_psf\_star} table. For magnitudes that are used to calculate colour, we use band-calibrated magnitudes from the \texttt{object} tables.
